\chapter{Nonmatching State Transfer and Restart}
\label{ch:transfer}

\section{Why state transfer is required}
Pre-refinement does not require state transfer because the final fracture calculation starts from the virgin state. Evolving remeshing is different: when the mesh changes after damage has developed, the receiving discretization must recover the mechanical displacement, phase field and irreversible history state of the donor analysis. A technically valid interpolation is not sufficient; the restarted model must also remain mechanically and thermodynamically consistent.

\section{Transferred quantities and mapping spaces}
The implemented transfer path distinguishes nodal and integration-point quantities. Displacement and phase field are transferred in nodal space using donor-element search and isoparametric interpolation. The crack-driving history field is stored at integration points and is therefore transferred with a separate Gauss-point mapping procedure. Crack flanks are treated as separate topological regions to avoid interpolation across an existing slit.

The receiving state is checked for phase bounds and history positivity. For irreversible fracture, preservation of the donor history is more important than reproducing a visually smooth field because a reduced history value can permit unphysical damage recovery.

\section{Runtime state ingestion}
An early transfer attempt generated a mapping artifact correctly but did not alter the running solution because the active UEL initialization path did not consume the prepared state. This distinction led to a redesign in which the transfer data are loaded explicitly at runtime through \texttt{UEXTERNALDB} and shared with the user-element routines. Subsequent tests confirmed that the transferred displacement, phase and history data were actually installed into the receiver analysis.

Figure~\ref{fig:transfer_phase_history} shows the actual checkpoint phase and history fields used by the transfer workflow. It complements the numerical transfer tables by making the donor localization state visible.

\begin{figure}[htbp]
\centering
\includegraphics[width=0.94\textwidth]{fig07_state_transfer_phase_history.png}
\caption{Checkpoint phase-field and history-field evidence used by the nonmatching state-transfer workflow. The panels originate from extracted simulation fields at the same donor checkpoint.}
\label{fig:transfer_phase_history}
\end{figure}

\section{Restart sequence}
To separate state installation from subsequent equilibration, the receiver calculation is divided conceptually into four stages:
\begin{enumerate}
  \item install the transferred state and prescribed donor kinematics;
  \item establish mechanical equilibrium while the phase field is held fixed;
  \item release the phase-field constraints and allow the coupled state to relax;
  \item continue the external loading.
\end{enumerate}
This sequence improved numerical robustness, but it must not be interpreted as a guarantee of physical continuity. The decisive comparison is the state after the phase field has been released and relaxed.

\begin{figure}[htbp]
\centering
\includegraphics[width=0.94\textwidth]{fig08_state_transfer_continuation.png}
\caption{Topology-changing transfer and continuation evidence across state installation, equilibration, phase release and continuation. The figure visualizes why successful ingestion and solver continuation do not by themselves establish restart equivalence.}
\label{fig:transfer_continuation}
\end{figure}

\section{Topology-changing restart results}
Two topology-changing Mode-II restart calculations reached the continuation stage but failed the predeclared reaction-force continuity check after phase release. A transfer based on an earlier donor state exhibited a force reduction of approximately $26.31\%$. A later donor state with a more strongly degraded crack-tip facet exhibited a reduction of approximately $99.56\%$.

\begin{figure}[htbp]
\centering
\includegraphics[width=0.78\textwidth]{fig_transfer_force_drop.pdf}
\caption{Magnitude of the reaction-force change between the mechanically equilibrated receiver state and the relaxed phase-release state for two tested topology-changing restarts. The results demonstrate that successful solver continuation is not equivalent to mechanically continuous state transfer.}
\label{fig:transfer_drop}
\end{figure}

The later case is particularly informative: the geometric facet selected for separation carried only a small local traction before the split, yet releasing the phase field caused a much larger structural response. This shows that a local crack-tip criterion alone is insufficient to establish restart equivalence; the receiving phase-field problem can reorganize over a larger process zone.

\section{Energy-accounting limitation}
The current UEL includes the Abaqus \texttt{ENERGY(8)} argument but does not populate a qualified internal-energy contribution. Consequently, global Abaqus variables such as \texttt{ALLSE} and \texttt{ALLIE} do not provide a complete UEL-aware fracture-energy balance for these restart comparisons. External work or $\tfrac12 Fu$ estimates are therefore not substituted for the missing internal-energy continuity check.

\section{Current conclusion}
The work demonstrates nonmatching data mapping and actual runtime ingestion of the transferred variables. It does not yet demonstrate a topology-changing restart that preserves the relaxed mechanical state. This distinction is retained in the remainder of the thesis: state-transfer infrastructure is available, while mechanically consistent evolving remeshing remains an open problem.
